%!TEX root = ../main.tex

In this section, we briefly discuss how to extend our results to general loss
functions.
Although we use the least-squares error throughout our derivations, this can be
generalized to a smooth and strictly convex loss function \( \calL : \R^n \times \R^n \into \R \)
without difficulty.
To do so, consider the more general problem,
\begin{equation}\label{eq:general-gl}
	\begin{aligned}
		p^*(\lambda) \!=\! \min_{w} \,
		 & F_\lambda(w) \!:=\! \half \calL(X w, y) + \lambda \sum_{\bi \in \calB} \norm{\wi}_2 \\
		 & \quad \text{s.t.} \quad \Ki^\top \wi \leq 0  \text{ for all } \bi \in \calB.
	\end{aligned}
\end{equation}
If \( \calL \) is strictly convex, then uniqueness of the
optimal model fit \( \hat y(\lambda) = X w \) follows from straightforward
adaption of \cref{lemma:unique-fit-cgl}.
Indeed, the only property of the squared-error used in this lemma is strict
convexity.

Since the model fit is constant in \( \solfn \) and \( \calL \) is both smooth and strictly convex,
the gradient \( \nabla_w \calL( X w, y) \), which is given by
\[
	X^\top \nabla_{\hat y} \calL( \hat y, y),
\]
must also be constant over \( \solfn \).
Thus, it is straightforward to replace the correlation vector \( \ci = \Xbi^\top (y - X w) \) with \( \Xbi^\top \nabla_{\hat y} \calL(\hat y, y) \)
throughout our derivations.

We form a Lagrange dual problem for CGL for one continuity-type result.
\cref{prop:model-fit-continuity} uses the Lagrange dual to
show that the correlation vector \( \ci \) is the unique solution to a convex
optimization program and applies standard sensitivity results to obtain continuity
of \( \hat y \).
In this same fashion, \( \Xbi^\top \nabla_{\hat y} \calL(\hat y, y) \) is the
unique solution to a Lagrange dual problem where the dual objective uses
the convex conjugate \( \calL^* \), rather than the dual of the quadratic penalty.
If \( \nabla_{\hat y} \calL(\hat y, y) \) is continuous in \( \hat y \), then
this is sufficient to deduce continuity of the model fit using the same
argument.
